\documentclass[12pt, a4paper]{article}

% --- Pakete für Sprache und Schriftart ---
\usepackage[utf8]{inputenc} % Ermöglicht Umlaute direkt im Text
\usepackage[ngerman]{babel} % Deutsche Silbentrennung und Bezeichnungen
\usepackage[T1]{fontenc}    % Schöne Schriftdarstellung
\usepackage{lmodern}        % Nutzt die modernere "Latin Modern" Schrift

% --- Layout & Design ---
\usepackage{geometry}
\geometry{left=2.5cm, right=2.5cm, top=3cm, bottom=3cm} % Seitenränder
\usepackage{microtype}     % Verbessert den Textsatz (weniger Trennungsfehler)

% --- Grafiken & Mathe ---
\usepackage{graphicx}      % Zum Einbinden von Bildern
\usepackage{amsmath}       % Mathematische Formeln
\usepackage{amssymb}

% --- Meta-Daten ---
\title{BWL - Klausur}
\author{Rebecca Schäfer}
\date{\today}

\begin{document}

\maketitle

\tableofcontents
\newpage

\section{SWOT - Analyse}

\subsection{Einleitung}
Die SWOT-Analyse ist ein Instrument zur Gegenüberstellung von innerbetrieblichen 
Stärken (Strengths) und Schwächen (Weaknesses) sowie externen Chancen 
(Opportunities) und Gefahren (Threats). Das Ziel ist es, aus dieser Analyse 
begründete strategische Empfehlungen für das Unternehmen abzuleiten. 
Während Stärken und Schwächen durch das Management beeinflussbar sind, 
entziehen sich Chancen und Gefahren der direkten Kontrolle. 
Gemäß den Buchkorrekturen sollte der Begriff „Threats“ im Deutschen besser 
mit „Gefahren“ oder „Bedrohungen“ übersetzt werden.

\subsection{Hauptteil: Analyse der Speed Motorradfabrik GmbH}
Basierend auf der Unternehmensbeschreibung (Situation 9) und den SWOT-Kopien lässt sich folgendes Profil erstellen:
Stärken (Strengths):\newline
\begin{itemize}
    \item Produktqualität: Absolute Spitzenqualität im Premiumsegment mit der Möglichkeit zu Spezialanfertigungen.
    \item Produktion: Moderne Anlagen und ein besonderes Fertigungskonzept („Einzelarbeit“: ein Mitarbeiter fertigt ein Motorrad von A bis Z).
    \item Personal: Hohe Motivation, sehr gutes Betriebsklima und ein kooperativer Führungsstil.
    \item Marketing: Hohe Weiterempfehlungsrate %(40 % Neukunden durch Mund-zu-Mund-Propaganda).
    \item Rentabilität: Überdurchschnittliche Umsatzrendite von 12,5 % in der Vergangenheit
\end{itemize}

Schwächen (Weaknesses):\newline
\begin{itemize}
    \item Kostenstruktur: Steigende Kostensituation und ein leichter Gewinnrückgang in den letzten Jahren.
    \item Personalfixierung: Mit nur 21 gewerblichen Mitarbeitern und Einzelfertigung könnte der Ausfall einzelner Fachkräfte die Produktion stark bremsen.
    \item Abhängigkeit: Fokus auf nur 3 Motorradtypen im Premiumsegment
\end{itemize}

Externe Analyse (Chancen & Gefahren)
Chancen (Opportunities):
\begin{itemize}
    \item Expansion: Der Inhaber plant aktiv eine regionale Expansion und weiteres Wachstum.
    \item Markterschließung: Ausbau der Internet-Präsenz über die bisherige Web-Seite hinaus, um neue Kundengruppen bundesweit anzusprechen.
\end{itemize}

Gefahren (Threats/Bedrohungen):
\begin{itemize}
    \item Marktentwicklung: Mögliche Änderungen der Kundenbedürfnisse (z. B. Trend zu Elektroantrieben, da die aktuellen Modelle „stark motorisiert“ sind).
    \item Wettbewerb: Zunehmender Konkurrenzdruck durch andere Premiumhersteller.
    \item Beschaffung: Preissteigerungen bei den benötigten hochwertigen Materialien und Komponenten.
\end{itemize}

\subsection{Strategische Ableitung (Vorschläge & Fazit)}
Durch die Kombination der Faktoren lassen sich Strategien ableiten:
\begin{itemize}
    \item S-O-Strategie (Ausbauen): Die Stärke der Einzelfertigung und Spitzenqualität nutzen, um in neue Märkte zu expandieren. Kunden schätzen das „Handmade in Ulm“-Gefühl, was bei einer Expansion als Alleinstellungsmerkmal dienen kann.
    \item W-O-Strategie (Aufholen): Um die steigenden Kosten (Schwäche) in den Griff zu bekommen und das geplante Wachstum (Chance) zu realisieren, könnten Prozesse in der Materialbeschaffung optimiert werden, ohne die Qualität zu gefährden.
    \item S-T-Strategie (Absichern): Den hohen Bekanntheitsgrad und die Mund-zu-Mund-Propaganda nutzen, um sich gegen neue Wettbewerber abzusichern (Image-Werbekampagne).
\end{itemize}
Fazit als Assistent der Geschäftsleitung: Die Speed Motorradfabrik ist gesund, muss aber die Kostenentwicklung 
genau beobachten, damit der Gewinnrückgang das geplante Wachstum nicht gefährdet. Ich empfehle, die Marketingaktivitäten 
im Internet zu intensivieren, um die Abhängigkeit von der Mund-zu-Mund-Propaganda zu verringern und die Expansion aktiv voranzutreiben

\section{Produkt - Markt - Matrix}
\subsection{Einleitung: Was ist die Produkt-Markt-Matrix?}
Die von Harry Igor Ansoff entwickelte Matrix (auch Marktfeldstrategien genannt) ist ein Instrument 
des strategischen Managements zur Planung des Unternehmenswachstums. Sie stellt zwei Dimensionen 
gegenüber: Produkte (gegenwärtig/neu) und Märkte (gegenwärtig/neu). Daraus ergeben sich vier 
grundlegende Wachstumsstrategien, um die „Zielelücke“ zwischen der erwarteten und der gewünschten Umsatzentwicklung zu schließen.
\subsection{Hauptteil: Anwendung auf die Speed Motorradfabrik GmbH}
In der Industrie verwenden wir gemäß den Buchkorrekturen nicht den Begriff Sortiment, sondern sprechen von unserem Leistungsprogramm, das aktuell aus drei Typen von Premium-Motorrädern sowie Spezialanfertigungen besteht.
\begin{itemize}
\item Marktdurchdringung (Gegenwärtige Produkte – Gegenwärtige Märkte): \newline
    Ziel: Erhöhung des Marktanteils mit dem bestehenden Leistungsprogramm im aktuellen Markt (Deutschland). \newline
    Maßnahme: Intensivierung der Werbung in Motorrad-Fachzeitschriften und gezielte 
    Förderung der Mund-zu-Mund-Propaganda, die bereits 40 Prozent der Neukunden generiert. 
    Ziel ist es, Kunden von Wettbewerbern abzuziehen oder Gelegenheitskäufer zu binden.
\item Marktentwicklung (Gegenwärtige Produkte – Neue Märkte):
    Ziel: Erschließung neuer geografischer Regionen oder Zielgruppen mit den vorhandenen Modellen.
    Maßnahme: Da Herr Heumann regional expandieren möchte, sollten wir über die bestehende Web-Seite hinaus das Internetmarketing professionalisieren, um Kunden im europäischen Ausland anzusprechen.
\item Produktentwicklung (Neue Produkte – Gegenwärtige Märkte): \newline
    Ziel: Den bestehenden Kunden völlig neuartige Produkte oder Produktvarianten anbieten. \newline
    Maßnahme: Entwicklung eines Elektro-Motorrads im Premiumsegment, um auf ökologische 
    Trends zu reagieren, oder Aufnahme einer exklusiven Fahrerausstattung (Bekleidung) in das Programm.
\item Diversifikation (Neue Produkte – Neue Märkte): \newline
    Ziel: Wachstum durch völlig neue Leistungen in neuen Märkten (höchstes Risiko).\newline
    Maßnahme: Beispielsweise eine laterale Diversifikation, indem wir unsere Expertise 
    in der hochpräzisen Einzelfertigung nutzen, um Komponenten für die Luftfahrt (Aluminiumherstellung) zu produzieren.
\end{itemize}

\subsection{Vorschläge und Fazit}
Als Ihr Assistent empfehle ich, den Fokus primär auf die Marktentwicklung zu legen, da dies dem ausdrücklichen Wunsch der Geschäftsleitung nach Expansion entspricht.
Strategische Überlegungen:
\begin{itemize}
\item Rentabilität: Unsere aktuelle Umsatzrendite von 12,5 Prozent ist exzellent und liegt deutlich über dem Branchenziel von 5–8 Prozent. 
    Wir müssen jedoch die steigende Kostensituation im Auge behalten, damit das Wachstum nicht zu Lasten des Gewinns geht.
\item Kapazitäten: Da wir in Einzelarbeit (ein Mitarbeiter pro Motorrad) fertigen, müssen wir bei einer Marktausweitung 
    rechtzeitig neue hochqualifizierte Fachkräfte einstellen, um die „Absolute Spitzenqualität“ zu sichern.
\item Risikomanagement: Die Marktdurchdringung ist die sicherste, die Diversifikation die risikoreichste Strategie. 
    Ein gesunder Mix aus Marktdurchdringung und regionaler Marktentwicklung ist für unsere Manufaktur am realistischsten.
\end{itemize}

\section{Produktlebenszyklus}
\subsection{Einleitung: Das Modell des Produktlebenszyklus}
Das Lebenszykluskonzept ist als Theorie oder Modell zu verstehen, das den Weg eines 
Produkts von der Markteinführung bis zum Ausscheiden aus dem Markt beschreibt. 
Es dient uns als strategisches Orientierungsinstrument, um den Erfolg 
unserer Premium-Motorräder über die Zeit (Umsatz- und Gewinnentwicklung) zu bewerten.

\subsection{Hauptteil: Der Produktlebenszyklus am Beispiel eines neuen Premium-Motorrads}

\subsubsection{A. Zeichnerische Beschreibung (Modell)}
Das Modell wird typischerweise in einem Koordinatensystem dargestellt (X-Achse: Zeit; Y-Achse: Umsatz/Gewinn),. Wir unterscheiden dabei fünf Phasen:
\begin{enumerate}
    \item Einführungsphase: Markteintritt.
    \item Wachstumsphase: Starker Anstieg von Umsatz und Gewinn.
    \item Reifephase: Umsatz erreicht Maximum, Gewinn beginnt bereits leicht zu sinken
    \item Sättigungsphase: Markt ist gesättigt, Preiskampf nimmt zu.
    \item Degenerationsphase (Rückgang): Umsatz und Gewinn fallen stark ab.
\end{enumerate}
   
Wichtige Anmerkung aus der Praxis: Gemäß den Buchkorrekturen wird die Gewinnkurve
in theoretischen Modellen oft viel zu hoch dargestellt, was nicht der wirtschaftlichen Realität entspricht.

\subsection{B. Die Phasen im Überblick}
Wenn die Speed Motorradfabrik ein neues Modell (z. B. eine innovative Spezialanfertigung) einführt, ist diese Phase durch folgende Merkmale gekennzeichnet:
\begin{itemize}
    \item Umsatz: Die Erlöse sind zunächst gering, da das Motorrad erst bekannt gemacht werden muss,.
    \item Gewinn: In dieser Phase erzielen wir in der Regel keine Gewinne, sondern Verluste. Ursache sind die hohen Kosten für die Produktentwicklung sowie hohe Marketingausgaben (z. B. für Anzeigen in Fachzeitschriften), um das Bekanntheitsgrad zu steigern,,.
    \item Produktion: Da wir in Einzelarbeit fertigen (ein Mitarbeiter pro Motorrad), sind die Stückkosten zu Beginn besonders hoch, da Lernprozesse erst einsetzen müssen.
    \item Ziel: Wir müssen „Meinungsführer“ und „Trendsetter“ im Premiumsegment überzeugen, damit das Produkt erfolgreich in die Wachstumsphase übergeht.
\ȩnd{itemize}
\subsubsection{C. Detailbetrachtung: Die Einführungsphase}

Pro (Vorteile) \newline
\begin{itemize}
    \item Anschaulichkeit: Zeigt sehr bildhaft das „Werden und Vergehen“ von Produkten
    \item Strategiehilfe: Erlaubt die Ableitung spezifischer Marketingaktivitäten (z. B. Erhöhung der Werbekosten in der Einführung)
    \item Frühwarnsystem: Hilft dabei, rechtzeitig Nachfolgeprodukte zu planen, bevor ein Modell in die Sättigung gerät.
\end{itemize}

Kontra (Nachteile/Kritik) \newline
\begin{itemize}
    \item Bestimmung der Phasen: Es ist in der Praxis extrem schwierig, exakt festzustellen, in welcher Phase man sich befindet.
    \item Keine Prognosekraft: Es lassen sich keine sicheren Aussagen über die tatsächliche Dauer oder die Höhe der Gewinne machen.
    \item Unrealistische Darstellung: Theoretische Gewinnkurven sind oft praxisfern und zu optimistisch.
\end{itemize}

\subsection{Fazit und Empfehlung}
Für die Speed Motorradfabrik ist das Modell wertvoll, um die Expansion von Herrn Heumann zu planen. 
Da wir aktuell eine sehr gute Umsatzrendite von 12,5 Prozent haben, verfügen wir über die nötigen finanziellen 
Mittel, um neue Modelle durch die verlustreiche Einführungsphase zu tragen. Wir sollten unsere starke 
Mund-zu-Mund-Propaganda (40 Prozent der Neukunden) nutzen, um die Einführungsphase neuer Modelle zu verkürzen 
und schneller die Gewinnzone (Break-even-Point) zu erreichen.

\section{BCG-Portfolio-Matrix}
\subsection{Einleitung: Was ist die BCG-Portfolio-Matrix?}
Die von der Boston Consulting Group entwickelte Matrix ist ein Instrument der Portfolio-Analyse, das strategische Geschäftseinheiten (SGE) anhand von zwei Dimensionen bewertet: dem Marktwachstum (Umweltkomponente) und dem relativen Marktanteil (Unternehmenskomponente). Ziel ist ein ausgewogenes Verhältnis zwischen traditionellen „Geldgebern“ und zukunftsorientierten Geschäftsfeldern.

\subsection{Hauptteil: Die vier Quadranten und Anwendung auf Speed Motorrad}
In der Matrix werden die Produkte als Kreise dargestellt, deren Größe dem jeweiligen Umsatz entspricht.
\begin{itemize}
\item Question Marks (Nachwuchsprodukte):
    \begin{itemize}
        \item Eigenschaften: Hohes Marktwachstum, aber noch niedriger Marktanteil (Einführungsphase).
        \item Situation Speed: Hierzu zählen unsere neu geplanten Modelle für die regionale Expansion oder innovative Prototypen mit alternativen Antrieben.
        \item Normstrategie: Offensivstrategie. Es muss entschieden werden, ob massiv investiert wird, um zum „Star“ zu werden, oder ob das Projekt abgebrochen wird.
    \end{itemize}
\item Stars (Sternprodukte):
    \begin{itemize}
        \item Eigenschaften: Hohes Marktwachstum und hoher Marktanteil (Wachstumsphase).
        \item Situation Speed: Unser aktuelles Top-Modell, das durch die rege Mund-zu-Mund-Propaganda (40 Prozent Neukunden) stark wächst.
        \item Normstrategie: Investitionsstrategie, um die Wettbewerbsposition zu halten oder auszubauen.
    \end{itemize}
\item Cash Cows (Melkkühe):
    \begin{itemize}
        \item Eigenschaften: Niedriges Marktwachstum, aber hoher Marktanteil (Reifephase).
        \item Situation Speed: Unsere etablierten 3 Standardtypen in Spitzenqualität, die die exzellente Umsatzrendite von 12,5 % erwirtschaften.
        \item Normstrategie: Abschöpfungsstrategie. Den Marktanteil halten und den Finanzmittelüberschuss nutzen, um die „Question Marks“ zu finanzieren.
    \end{itemize}
\item Poor Dogs (Auslaufprodukte):
    \begin{itemize}
        \item Eigenschaften: Niedriges Marktwachstum und niedriger Marktanteil (Degenerationsphase).
        \item Situation Speed: Ältere Spezialanfertigungen oder Modelle, deren Nachfrage aufgrund des Marktwandels stagniert.
        \item Normstrategie: Desinvestitionsstrategie. Diese Produkte sollten eliminiert oder verkauft werden, um Ressourcen freizusetzen.
    \end{itemize}
\end{itemize}
\subsection{Vorschläge und Fazit}
\begin{itemize}
\item Investitionskraft nutzen: Die hohen Gewinne unserer Cash Cows müssen gezielt in die Question Marks fließen, 
    um das Ziel von Herrn Heumann – weiteres Wachstum – zu erreichen.
\item Kostenkontrolle: Trotz der guten Marktposition müssen wir die steigende Kostensituation im 
    Auge behalten. In der Theorie wird die Gewinnkurve oft unrealistisch hoch dargestellt; wir müssen sicherstellen, dass unsere reale Rentabilität stabil bleibt.
\item Kapazitätsmanagement: Da wir in Einzelarbeit fertigen, müssen wir bei einem Aufstieg von 
    „Question Marks“ zu „Stars“ rechtzeitig unsere personellen Kapazitäten (derzeit 21 gewerbliche Mitarbeiter) planen, damit die Spitzenqualität nicht leidet.
\end{itemize}



\section{Bilder einbinden}
Um Bilder zu nutzen, kannst du sie einfach in einen Unterordner (z.B. \texttt{images/}) legen und so einbinden:
% \begin{figure}[h]
%     \centering
%     \includegraphics[width=0.5\textwidth]{images/mein-bild.png}
%     \caption{Eine Bildbeschreibung}
%     \label{fig:mein-bild}
% \end{figure}

\end{document}